% Sempre deve haver um resumo em português. Parágrafo único, com no
% máximo 500 palavras e numa única página; palavras-chave logo abaixo,
% separadas por ponto e vírgula e finalizadas por ponto (NBR 6028).
\begin{resumo}
Exoesqueletos passivos de apoio ao braço são empregados no trabalho com os braços elevados, associado a distúrbios do ombro, e o auxílio que entregam depende apenas da geometria do mecanismo e da mola. As patentes que descrevem os modelos comerciais detalham o mecanismo, mas não informam massa, inércia nem rigidez. Este trabalho reconstrói, a partir da documentação de patente, três mecanismos de arquitetura distinta e formula modelos estáticos do torque que eles entregam em função do ângulo de elevação, confrontando-o com o torque gravitacional do braço ou com a curva publicada: o gerador de torque de mola, cabo e polia da US 9.889.554 B2, o dispositivo da EP 3 504 033 B1, em que um elemento elástico atua diretamente sobre uma alavanca solidária ao braço, e o mecanismo da US 9.737.374 B2, em que a mola age por um came de raio variável. A geometria foi reconstruída em CAD, a cinemática foi resolvida em forma fechada, e cada etapa do cálculo, implementado em Python, foi conferida por verificações internas de consistência, entre elas o princípio dos trabalhos virtuais, que mostrou, na US 9.889.554 B2, que a extensão da mola deve acompanhar o caminho completo do cabo. Com os parâmetros não publicados congelados, dois deles de origem não reproduzida, o modelo do primeiro mecanismo reproduz a forma da curva publicada, nos mesmos 29 pontos usados para fixá-los, com coeficiente de determinação de 0,99839 e raiz do erro quadrático médio de 1,24\% do maior valor extraído, e o pico do torque do cabo, a 97,5° de elevação, fica cerca de 8° acima do pico gravitacional, de 11,16 N·m a 89,4°. Como a ordenada da curva não tem unidade declarada, o modelo determina a forma do torque, mas não sua escala física. Na EP 3 504 033 B1, o espaçamento entre as três curvas de pré-tensão publicadas determina os comprimentos de referência do elástico, de 300,91 mm, 287,19 mm e 275,59 mm, sem depender da escala da ordenada nem da rigidez. A forma prevista com a geometria do CAD, contudo, não reproduz as curvas da patente, e a rigidez não é identificável pela figura; sob a hipótese de assistência de 40\% do torque do peso em 90°, o pico do torque é de 4,376 N·m em 78,8°. Nesses dois dispositivos, o pico de assistência fica afastado do pico gravitacional, acima dele no primeiro e abaixo no segundo, e a comparação entre eles fica respondida quanto aos ângulos e à forma dos perfis, mas não quanto à magnitude. Na US 9.737.374 B2, com o perfil do came extraído do desenho, uma mola linear reproduz a curva publicada com coeficiente de determinação de 0,901, ou de 0,953 com o came girado 49°, giro que também alinha as zonas do came às da patente. As etapas dinâmicas, pela formulação de Lagrange e no MBDyn, dependem de um dado independente que fixe a escala do torque e não foram executadas.

\bigskip\noindent Palavras-chave: exoesqueleto passivo; compensação de gravidade; ombro; patentes; identificação de parâmetros.
\end{resumo}

% Sempre deve haver um abstract, com as mesmas regras.
\begin{abstract}
Passive arm-support exoskeletons are used in work with raised arms, which is associated with shoulder disorders, and the assistance they deliver depends only on the geometry of the mechanism and on the spring. The patents that describe the commercial models detail the mechanism but do not report mass, inertia or stiffness. This work reconstructs, from patent documentation, three mechanisms of different architecture and formulates static models of the torque they deliver as a function of the elevation angle, comparing it with the gravitational torque of the arm or with the published curve: the spring, cable and pulley torque generator of US 9,889,554 B2, the device of EP 3 504 033 B1, in which an elastic element acts directly on a lever attached to the arm, and the mechanism of US 9,737,374 B2, in which the spring acts through a variable-radius cam. The geometry was reconstructed in CAD, the kinematics was solved in closed form, and each step of the calculation, implemented in Python, was verified by internal consistency checks, among them the principle of virtual work, which showed, for US 9,889,554 B2, that the spring extension must follow the complete cable path. With the unpublished parameters frozen, two of them of unreproducible origin, the model of the first mechanism reproduces the shape of the published curve, on the same 29 points used to fix them, with a coefficient of determination of 0.99839 and a root mean square error of 1.24\% of the largest extracted value, and the peak of the cable torque, at 97.5° of elevation, lies about 8° above the gravitational peak of 11.16 N·m at 89.4°. Since the ordinate of the curve has no declared unit, the model determines the shape of the torque but not its physical scale. For EP 3 504 033 B1, the spacing between the three published pretension curves determines the reference lengths of the elastic element, 300.91 mm, 287.19 mm and 275.59 mm, without depending on the ordinate scale or on the stiffness. The shape predicted with the CAD geometry, however, does not reproduce the patent curves, and the stiffness cannot be identified from the figure; under the assumption of assistance equal to 40\% of the weight torque at 90°, the torque peak is 4.376 N·m at 78.8°. In these two devices, the assistance peak is set apart from the gravitational peak, above it in the first and below it in the second, and the comparison between them is answered regarding the angles and the shape of the profiles, but not regarding magnitude. For US 9,737,374 B2, with the cam profile extracted from the drawing, a linear spring reproduces the published curve with a coefficient of determination of 0.901, or 0.953 with the cam turned 49°, a rotation that also aligns the cam zones with those of the patent. The dynamic steps, through the Lagrange formulation and in MBDyn, depend on independent data that fix the torque scale and were not carried out.

\bigskip\noindent Keywords: passive exoskeleton; gravity compensation; shoulder; patents; parameter identification.
\end{abstract}

% Se houver um resumo em espanhol, descomente as linhas a seguir:
%\begin{resumen}
% A mesma regra aplica-se.
%\end{resumen}
